\answer{ex:07:1}{$V=\kappa p/\bigl(\kappa p+(1-\kappa)(1-p)\bigr)$：$0.059$、$0.20$、$0.50$；中位数为$0$，而点~\eqref{eq:expectile}随$\kappa$平滑变化。}
\answer{ex:07:2}{$p=1/3$，$x=0.75$，$\sigma_r=0.35$，$V(0.2)=0.083$。}
\answer{ex:07:3}{$V=\sqrt\kappa/(\sqrt\kappa+\sqrt{1-\kappa})$，从$0.25$到$0.75$；离散$\propto\sqrt\eta$；两滴情形下$V=0.25+0.5\kappa$，两种分布在两端的神经元上相差$0.05$，需要$\eta\lesssim0.01$。}
\answer{ex:07:4}{（a）~$J^*=2\sqrt{C\bar R}$。（b）~多付$\bigl(k^{1/2}+k^{-1/2}\bigr)/2-1$：$k=2$时为$6\,\%$，$k=4$时为$25\,\%$——“两倍以内”的精度就够了。}
\answer{ex:07:5}{尖峰面积为$R(e^{-1}-e^{-2})\approx0.23R$——相当于$2.3\,\text{s}$的上升；如果多巴胺报告$V$，就没有尖峰，电平只是台阶式地增大为$e$倍。}
\answer{ex:07:6}{事件使权重移动$\propto A\tau_f/(\tau_e+\tau_f)$，电平带来误差$s\tau_f$：$9\,\text{s}\le\tau_f\le17\,\text{s}$。}
\answer{ex:07:7}{初始时$\Delta P=0.80$，$M=0.90$。（a）~$\Delta P=0.74$（$-8\,\%$），$M=0.43$（$-52\,\%$）；（b）~$\Delta P=0.40$（$-50\,\%$），$M=0.80$（$-11\,\%$）：双重分离。}
\answer{ex:07:8}{$N_{\text{eff}}\to2+\rho^2N\approx10^6$次尝试：比单根导线少$10^4$倍，但$1\,\%$的泄漏仍要花掉一百万次尝试。}
